\section{From transverse pulses to fluid motion}
\label{sv:imports}

\subsection{Transverse pulses and classical uniqueness}
\label{sv:analytic}
A transverse field has either one nonzero component independent of its
direction of motion, or the form $(a(z),b(z),0)$.  Both its divergence
and its self-advection vanish.  A finite sum of such fields need not have
zero self-advection if several act simultaneously; we instead give them
disjoint time supports.  This scheduling restriction is part of every
construction below.

\begin{lemma}[Direct solenoidal forcing]\label{sv:force}
Let $U(t,x)$ on a flat torus be a smooth sequence of transverse pulses
with disjoint temporal supports, each of zero spatial mean, and let
$U(0)=0$.  Assume $U$ and its derivatives are bounded on each finite
time interval.  Then $f=U_t-\nu\Delta U$ is smooth, solenoidal and
mean zero, and $(U,0)$ solves \eqref{sv:NS}.  Its velocity and normalized
pressure are unique in the classical comparison class defined above.
Its material flow exists for every finite forward time and is a
volume-preserving smooth diffeomorphism.  If a finite pulse sequence is
repeated, with zero collars at its joins, all mixed derivatives are
globally bounded and the energy is uniformly bounded.  These conclusions
also hold when one separate finite loading sequence is prepended.
\end{lemma}
\begin{proof}
At each time at most one transverse field is active, so
$\operatorname{div}U=0$ and $(U\cdot\nabla)U=0$.
Divergence commutes with $\partial_t$ and $\Delta$, and a periodic
Laplacian has zero integral.  These facts prove the force assertions and
the equation by substitution.  For another classical solution $v$, put
$w=v-U$.  Subtract the equations, multiply by $w$, and integrate on
the torus.  Incompressibility and periodic integration by parts cancel
the pressure and transport by $v$, giving
\[
 \frac12\frac d{dt}\|w\|_2^2+\nu\|\nabla w\|_2^2
 \le\|\nabla U\|_{\infty,\mathrm{op}}\|w\|_2^2.
\]
The stated classical regularity justifies each operation.  The coefficient
is bounded on every finite interval and $w(0)=0$, so an integrating
factor gives $w=0$.  The pressure gradient then vanishes, and its
normalization makes the pressure zero.  This is a comparison for the
displayed global solution, not an existence assertion for general data.

The spatial periodic lift of $U$ is bounded and Lipschitz on every
finite interval.  Local ODE existence, uniqueness and continuation give
its trajectories, and backward integration supplies their inverses.
Smooth dependence gives smooth diffeomorphisms.  The Jacobian obeys
$\partial_t\det D_aX=(\operatorname{div}U)\circ X\det D_aX$
with initial value one, proving volume preservation.  A finite loading
interval followed by a fixed smooth repeated pattern has a compact
collection of phase templates; their derivative bounds imply all the
last assertions.
\end{proof}

\subsection{The reciprocal geometric input}
Theorem~3.2 of~\cite{Entry} has the following precise interface.
Fix rational $L,h>0$ and rational coordinate charts on $\mathbb T_L^3$.
Take finite families of closed axis-parallel rational rectangles $P_i,Q_i$
with positive side lengths and centers $p_i,q_i$, each family pairwise
positively separated, and maps
$F_i(y)=q_i+\operatorname{diag}(r_i,r_i^{-1})(y-p_i)$ with
$F_i(P_i)=Q_i$ and $r_i>0$ rational.  Suppose every source and target
half-width is at most $h$, their centers lie in a rational rectangle
$K$, and $K+[-2h,2h]^2$ is compactly contained in the planar chart.
For any rational coding height $z_0$ inside its chart, that theorem supplies
seven effective mean-zero transverse pulses, periodic with zero collars,
whose period map is $(y,z)\mapsto(F_i(y),z)$ on a positive
effectively specified neighborhood of each $P_i\times\{z_0\}$.  For a point
starting on $P_i$ at height $z_0$, the four deformation stages
remain within sup distance $2h$ of the source center;
the following translation interpolates between centers with a fixed
offset of sup norm at most $h$.  Lifting and lowering preserve the planar
coordinates.  Source--target intersections are allowed.  The theorem
also proves the empty-list case and the force conclusions of
Lemma~\ref{sv:force}.

This seven-stage construction supplies the reciprocal template used below.
Alternative shear lists and translation routes are checked locally,
including their cutoff plateaus and intermediate paths.  The constructions
on boxes of prescribed thickness likewise prove their own bounds on every
point of those boxes.
